\documentclass[a4paper,10pt]{article}
\usepackage[utf8]{inputenc}
\usepackage[T1]{fontenc}
\usepackage[french]{babel}
\usepackage{amsmath,amssymb}
\usepackage{geometry}
\usepackage{array,tabularx}
\usepackage{tikz}
\geometry{margin=1.6cm}
\pagestyle{empty}
\newcommand{\ligne}{\par\vspace{0.55cm}\noindent\dotfill\par}
\newcommand{\carreaux}[1]{\begin{center}\begin{tikzpicture}\draw[step=0.5cm,gray!35,very thin] (0,0) grid (17,#1);\end{tikzpicture}\end{center}}
\newenvironment{exercice}[2]{\paragraph{Exercice #1 \hfill #2 points}\mbox{}\par}{\medskip}
\begin{document}
\begin{center}\Large\bfseries Devoir surveillé 11 - Proportionnalité 2\end{center}
\vspace{2mm}
\begin{exercice}{1}{4}Un vélo parcourt $18$ km en $45$ min. En supposant la vitesse constante, quelle distance parcourt-il en $1$ h ? en $2$ h $30$ ?\end{exercice}\begin{exercice}{2}{5}Un pantalon à $48$ euros bénéficie d'une remise de $25\%$. Calculer le montant de la remise puis le prix payé.\end{exercice}\begin{exercice}{3}{5}Compléter le tableau de proportionnalité.\begin{center}\begin{tabular}{|c|c|c|c|c|}\hline Masse (kg)&2&5&8&12\\\hline Prix (euros)&7&&&\\\hline\end{tabular}\end{center}\end{exercice}\begin{exercice}{4}{6}Une carte est à l'échelle $1/25000$. Deux villages sont séparés de $7,2$ cm sur la carte. Calculer la distance réelle en km.\end{exercice}\end{document}
